\subsubsection{文件打开与关闭}
在src文件夹的file.c文件中实现对文件的打开与关闭操作,由于这里每个进程实际都有可能是写者和读者,因此打开文件时不区分是读进程还是写进程,只根据打开文件的模式来判断,具体代码如下:
\begin{itemize}
    \item 打开文件
    \begin{ccode}{打开文件函数}
int fs_open(int pid, const char *filename, int mode)
{
    if (pid < 0 || pid >= MAX_PROCESS || filename == NULL) {
        return -1;
    }

    int cur_dir = ProcessTable[pid].current_dir;
    int dir_index = -1;

    //互斥访问Directory,查找要打开的文件
    pthread_mutex_lock(&fs_mutex);

    for (int i = 0; i < MAX_ENTRY; i++) {
        if (Directory[i].ac == 1 &&
            Directory[i].parent == cur_dir &&
            Directory[i].type == 0 &&
            strcmp(Directory[i].name, filename) == 0) {
            dir_index = i;
            break;
        }
    }

    if (dir_index == -1) {
        pthread_mutex_unlock(&fs_mutex);
        printf("open failed: file not found: %s\n", filename);
        return -1;
    }

    pthread_mutex_unlock(&fs_mutex);

    //只读、读写打开文件需要修改reader_count数量(互斥修改)

    if (mode == MODE_READ) {
        pthread_mutex_lock(&Directory[dir_index].mutex);//互斥修改reader_count
        Directory[dir_index].reader_count++;
        if (Directory[dir_index].reader_count == 1) {
            if (sem_trywait(&Directory[dir_index].rw_lock) != 0) { //有写者时打开失败
                Directory[dir_index].reader_count--;
                pthread_mutex_unlock(&Directory[dir_index].mutex);
                printf("open failed: file is busy\n");
                return -1;
            }
        }
        pthread_mutex_unlock(&Directory[dir_index].mutex);//修改reader_count结束
    }
    else if (mode == MODE_WRITE || mode == MODE_RDWR) {
        if (sem_trywait(&Directory[dir_index].rw_lock) != 0) {  //文件被读者或写者占用时访问失败
            printf("open failed: file is busy\n");
            return -1;
        }
    }
    else {
        printf("open failed: invalid mode\n");
        return -1;
    }

    //分配fd
    for (int fd = 0; fd < MAX_FD; fd++) {
        if (ProcessTable[pid].fd_table[fd].used == 0) {
            ProcessTable[pid].fd_table[fd].used = 1;
            ProcessTable[pid].fd_table[fd].dir_index = dir_index;
            ProcessTable[pid].fd_table[fd].offset = 0;
            ProcessTable[pid].fd_table[fd].mode = mode;

            printf("open success: pid=%d, file=%s, fd=%d\n",
                   pid, filename, fd);
            return fd;
        }
    }

    //fd分配失败,回退已经做的修改
    if (mode == MODE_READ) {
        pthread_mutex_lock(&Directory[dir_index].mutex);
        Directory[dir_index].reader_count--;
        if (Directory[dir_index].reader_count == 0) {
            sem_post(&Directory[dir_index].rw_lock);  // 最后一个读者释放文件 
        }
        pthread_mutex_unlock(&Directory[dir_index].mutex);
    } else {
        sem_post(&Directory[dir_index].rw_lock);  // 写者释放文件 
    }

    printf("open failed: fd table full\n");
    return -1;
}
    \end{ccode}
    整体实现思路类似任务一,在Directory数组中查找要打开的文件时前后加锁保护,确保一致性。之后按读者-写者问题的经典解法实现:
    当以只读模式(MODE\_READ)打开时,该进程视为读者,需要先互斥修改reader\_count(通过pthread\_mutex\_lock/unlock对mutex加锁/解锁),若为第一个读者,还需对资源信号量执行P操作(sem\_trywait(\&Directory[dir\_index].rw\_lock)),以阻塞后续写者;
    当以只写模式(MODE\_WRITE)或读写模式(MODE\_RDWR)打开时,该进程视为写者,直接对资源信号量执行P操作,确保对文件的互斥访问,阻止其他读者和写者。
    这里实现时使用\texttt{sem\_trywait}而非\texttt{sem\_wait},原因是若使用sem\_wait阻塞等待,当前交互界面会被卡住,无法手动切换进程并释放资源。因此本实验采用非阻塞方式,即文件被占用时直接返回失败。在实际多进程环境中应当使用真正的阻塞等待来实现互斥。
    对于打开的文件需要分配文件描述符fd,当fd表已满时意味着文件打开失败,此时需要回退前一步做的修改,即恢复reader\_count以及释放资源信号量(sem\_post)。
    \item 关闭文件
    \begin{ccode}{关闭文件函数}
bool fs_close(int pid, int fd)
{
    if (pid < 0 || pid >= MAX_PROCESS || fd < 0 || fd >= MAX_FD) {
        return false;
    }

    if (ProcessTable[pid].fd_table[fd].used == 0) {
        printf("close failed: invalid fd\n");
        return false;
    }

    int dir_index = ProcessTable[pid].fd_table[fd].dir_index;
    int mode = ProcessTable[pid].fd_table[fd].mode;


    if (mode == MODE_READ) {
        pthread_mutex_lock(&Directory[dir_index].mutex);
        Directory[dir_index].reader_count--;
        if (Directory[dir_index].reader_count == 0) {
            sem_post(&Directory[dir_index].rw_lock);  // 最后一个读者释放文件 
        }
        pthread_mutex_unlock(&Directory[dir_index].mutex);
    }
    else if (mode == MODE_WRITE || mode == MODE_RDWR) {
        sem_post(&Directory[dir_index].rw_lock);  // 写者释放文件 
    }
    
    //释放已分配的fd
    ProcessTable[pid].fd_table[fd].used = 0;
    ProcessTable[pid].fd_table[fd].dir_index = -1;
    ProcessTable[pid].fd_table[fd].offset = 0;
    ProcessTable[pid].fd_table[fd].mode = 0;

    printf("close success: pid=%d, fd=%d\n", pid, fd);
    return true;
}
    \end{ccode}
     同样采用读者-写者问题的经典解法,类似打开文件的实现。根据进程数组中对应文件描述符记录的打开模式确定是读者还是写者:
     写者(MODE\_WRITE或MODE\_RDWR)直接释放资源信号量即可;读者(MODE\_READ)需要先互斥减少reader\_count,当reader\_count降为0时由最后一个读者释放资源信号量。
\end{itemize}

\subsection{任务二函数汇总}
下表对任务二实现的所有主要函数进行汇总,以便对照查阅。

\begin{table}[H]
\centering
\caption{任务二主要函数汇总}
\label{tab:task2_summary}
\resizebox{\textwidth}{!}{%
\begin{tabular}{|c|c|c|c|c|}
\hline
\textbf{函数名} & \textbf{功能} & \textbf{实现方法} & \textbf{关键数据结构} & \textbf{同步机制} \\
\hline
\texttt{init\_fs()} & 初始化文件系统 & 新增信号量/互斥锁初始化,调用init\_process\_table() & Directory[].rw\_lock, .mutex & 信号量初值置1,互斥锁初始化 \\
\hline
\texttt{init\_process\_table()} & 初始化PCB数组 & memset清零,设置pid和current\_dir & ProcessTable[] & 无 \\
\hline
\texttt{mkdir(pid,...)} & 创建目录 & 互斥锁保护Directory数组,查找空闲项设置属性 & Directory[], fs\_mutex & 互斥锁保护全局目录项 \\
\hline
\texttt{ls(pid,...)} & 列出当前目录内容 & 互斥锁保护下遍历Directory并打印 & Directory[], fs\_mutex & 互斥锁保护读取一致性 \\
\hline
\texttt{cd(pid,...)} & 切换当前目录 & 互斥锁保护下查找目标目录,更新进程current\_dir & Directory[], ProcessTable[] & 互斥锁防止mkdir干扰 \\
\hline
\texttt{fs\_create(pid,...)} & 创建文件 & 互斥锁保护Directory,找空闲项设置属性 & Directory[], fs\_mutex & 互斥锁 \\
\hline
\texttt{fs\_delete()} & 删除文件 & 遍历Directory找文件,沿FAT链释放块 & Directory[], FAT[], fs\_mutex & 互斥锁 \\
\hline
\texttt{fs\_open(pid,...)} & 打开文件 & 互斥锁查找文件→读者:加reader\_count,首读者P(rw\_lock);写者:直接P(rw\_lock)→分配fd & Directory[].rw\_lock, .mutex & 读者-写者经典解法,sem\_trywait非阻塞 \\
\hline
\texttt{fs\_close(pid,...)} & 关闭文件 & 读者:减reader\_count,末读者V(rw\_lock);写者:V(rw\_lock)→清零fd表项 & Directory[].rw\_lock, .mutex & 读者-写者释放逻辑 \\
\hline
\texttt{fs\_read(pid,...)} & 读取文件内容 & 沿FAT链从offset逐块读取到buffer & Disk[], FAT[], ProcessTable[] & 依赖open时的rw\_lock保护 \\
\hline
\texttt{fs\_write(pid,...)} & 写入文件内容 & 确保FAT链有足够块,从offset逐块写入 & Disk[], FAT[], ProcessTable[] & 依赖open时的rw\_lock保护 \\
\hline
\end{tabular}%
}
\end{table}
